\par\Needspace{6\baselineskip}\medskip
\hypertarget{native-084-article065}{}\pdfbookmark[2]{65.}{article-065}
\begin{center}65.\end{center}
Sind aus gegebenen Elementen nur sehr wenige geozentrische Orte zu bestimmen, so lohnt es sich kaum, alle oben angegebenen Kunstgriffe anzuwenden, mit denen man von der exzentrischen Anomalie unmittelbar zur geozentrischen Länge und Breite oder sogar zur Rektaszension und Deklination übergehen kann. Die so erzielte Ersparnis würde durch die Vielzahl der zuvor zu berechnenden Hilfsgrößen aufgezehrt. Verbindet man aber die Reduktion auf die Ekliptik mit der Berechnung der geozentrischen Länge und Breite, so ergibt sich stets eine beachtliche Ersparnis. Wählt man nämlich die Ekliptik selbst als Ebene für die Koordinate $z$ und legt die Pole der Ebenen für die Koordinaten $x$, $y$ auf die Längen $\AscendingNode$, $90^0+\AscendingNode$, so lassen sich die Koordinaten sehr leicht ohne jede Hilfsgröße bestimmen. Man erhält
\SourceLeaf{085}
\[
\begin{array}{@{}l|l|l@{}}
x=r\cos u & X=R'\cos(L-\AscendingNode) & x-X=\Delta'\cos(l-\AscendingNode)\\[.4em]
y=r\cos i\sin u & Y=R'\sin(L-\AscendingNode) & y-Y=\Delta'\sin(l-\AscendingNode)\\[.4em]
z=r\sin i\sin u & Z=R'\operatorname{tang}B & z-Z=\Delta'\operatorname{tang}b.
\end{array}
\]
\noindent Ist $B=0$, so gilt $R'=R$, $Z=0$. Nach diesen Formeln wird unser Beispiel mit den folgenden Zahlen durchgerechnet. Ausgangspunkt ist $L-\AscendingNode=213^0\,12^{\prime}\,0^{\prime\prime}{,}32$
\begin{center}
\begin{tabular}{@{}ll@{}}
\makebox[13em][l]{$\log r$\hfill .\,.\,.\,.\,.\,.\,.\,.\,.\,.\,.} & $0{,}325\,9877$\\[.3em]
\makebox[13em][l]{$\log\cos u$\hfill .\,.\,.\,.\,.\,.\,.\,.} & $9{,}982\,4141_n$\\[.3em]
\makebox[13em][l]{$\log\sin u$\hfill .\,.\,.\,.\,.\,.\,.\,.} & $9{,}445\,4714_n$\\[.3em]
\hline
\makebox[13em][l]{$\log x$\hfill .\,.\,.\,.\,.\,.\,.\,.\,.\,.\,.} & $0{,}308\,4018_n$\\[.3em]
\hline
\makebox[13em][l]{$\log r\sin u$\hfill .\,.\,.\,.\,.\,.\,.} & $9{,}771\,4591_n$\\[.3em]
\makebox[13em][l]{$\log\cos i$\hfill .\,.\,.\,.\,.\,.\,.\,.} & $9{,}988\,5266$\\[.3em]
\makebox[13em][l]{$\log\sin i$\hfill .\,.\,.\,.\,.\,.\,.\,.} & $9{,}355\,7570$\\[.3em]
\hline
\makebox[13em][l]{$\log y$\hfill .\,.\,.\,.\,.\,.\,.\,.\,.\,.\,.} & $9{,}759\,9857_n$\\[.3em]
\makebox[13em][l]{$\log z$\hfill .\,.\,.\,.\,.\,.\,.\,.\,.\,.\,.} & $9{,}127\,2161_n$
\end{tabular}
\end{center}
\begin{center}
\begin{tabular}{@{}ll@{}}
\makebox[13em][l]{$\log R'$\hfill .\,.\,.\,.\,.\,.\,.\,.\,.\,.} & $9{,}998\,0979$\\[.3em]
\makebox[13em][l]{$\log\cos(L-\AscendingNode)$\hfill .\,.\,.} & $9{,}922\,6027_n$\\[.3em]
\makebox[13em][l]{$\log\sin(L-\AscendingNode)$\hfill .\,.\,.} & $9{,}738\,4353_n$\\[.3em]
\hline
\makebox[13em][l]{$\log X$\hfill .\,.\,.\,.\,.\,.\,.\,.\,.\,.} & $9{,}920\,7006_n$\\[.3em]
\makebox[13em][l]{$\log Y$\hfill .\,.\,.\,.\,.\,.\,.\,.\,.\,.} & $9{,}736\,5332_n$\\[.3em]
$Z=$ & $0$
\end{tabular}
\end{center}
\noindent Daraus ergibt sich
\begin{center}
\begin{tabular}{@{}ll@{}}
\makebox[13em][l]{$\log(x-X)$\hfill .\,.\,.\,.\,.\,.} & $0{,}079\,5906_n$\\[.3em]
\makebox[13em][l]{$\log(y-Y)$\hfill .\,.\,.\,.\,.\,.} & $8{,}480\,7165_n$\\[.3em]
woraus $(l-\AscendingNode)=$ & $181^0\,26^{\prime}\,33^{\prime\prime}{,}49$\\[.3em]
\makebox[13em][l]{$\log\Delta'$\hfill .\,.\,.\,.\,.\,.\,.\,.\,.\,.} & $0{,}079\,7283$\\[.3em]
\makebox[13em][l]{$\log\operatorname{tang}b$\hfill .\,.\,.\,.\,.\,.\,.} & $9{,}047\,4878_n$
\end{tabular}
\end{center}
\begin{center}
\begin{tabular}{@{}rl@{}}
$l=$ & $352^0\,34^{\prime}\,22^{\prime\prime}{,}22$\\[.4em]
$b=$ & $-6\quad21\quad55{,}06$
\end{tabular}
\end{center}
